% CN 5.16 Exercises (p.178--182)

% ======================================================================
\begin{exercise}{1}
    What can a predicate with no arguments be?
\end{exercise}

\begin{exercise}{2}
    Predicates: $\mathrm{prolog}(X)$: $X$ knows Prolog;
    $\mathrm{elvish}(X)$: $X$ knows Elvish.
    \begin{enumerate}[label=(\alph*)]
        \item Write a universal statement meaning ``Nobody knows both Prolog and
              Elvish.''
        \item Suppose (a) is False. Starting with its negation, derive an
              existential statement meaning that someone knows both languages.
    \end{enumerate}
\end{exercise}

\begin{exercise}{3}
    $\mathrm{supervised}(X,Y)$ means person $X$ supervised the PhD of person $Y$.
    Express each in predicate logic.
    \begin{enumerate}[label=(\alph*)]
        \item Alonzo Church was the PhD supervisor of Alan Turing.
        \item No person supervised their own PhD.
        \item Not everyone has supervised someone else's PhD.
        \item Not every PhD supervisor had a PhD supervisor. (Write this as an
              existential statement.)
        \item Someone had at least two PhD supervisors.
        \item Alan Turing supervised the PhDs of exactly two people.
    \end{enumerate}
\end{exercise}

\begin{exercise}{4}
    $\mathrm{designedFirstComputerIn}(X,Y)$ means person $X$ was one of the
    designers of the first computer in country $Y$. Express in predicate logic:
    \begin{enumerate}[label=(\alph*)]
        \item Trevor Pearcey and Maston Beard designed the first computer in
              Australia.
        \item No-one designed the first computer in two different countries.
    \end{enumerate}
\end{exercise}

\begin{exercise}{5}
    (Six degrees of separation) Consider the binary predicate
    $\mathrm{knows}$ (p.68, \S2.14) on the set of all people. It has been claimed
    that the distance between any two people in the social network is at most 6.
    Write this claim in predicate logic, using just $\mathrm{knows}$.
\end{exercise}

\begin{exercise}{6}
    Suppose you have the equality predicate for sets, the symmetric difference
    function for sets, and set variables $A$ and $B$. Write the statement of
    Theorem 4 (\emph{For any two sets $A$ and $B$, they are identical if and
    only if their symmetric difference is empty}) in predicate logic.
\end{exercise}

% ======================================================================
\begin{exercise}{7}
    For each statement, (i) identify the predicates, functions, variables (with a
    suggested domain for each) and constants used; (ii) state whether or not it
    is a proposition, and if so, whether it is true or false.
    \begin{enumerate}[label=(\alph*)]
        \item $\forall x\ x^2\ge0$
        \item $x^2<0$
        \item $\forall x\ 2^x=x^2$
        \item $x<0\Rightarrow\forall y\,((0<y)\Rightarrow(x<y))$
        \item $\forall x\forall z\exists y\ (x<y)\wedge(y<z)$
        \item $\exists x\exists y\ \neg(x\Rightarrow y)$
    \end{enumerate}
\end{exercise}

\begin{exercise}{8}
    Consider
    \[
        \forall A\forall B\ (A\subseteq B)\Rightarrow(\power A\subseteq\power B).
    \]
    What predicates, functions and variables does this statement use? Restate,
    in words, the theorem being stated.
\end{exercise}

\begin{exercise}{9}
    Let $W$ range over English words, $L$ over English letters, and
    $\mathrm{WordContainsLetter}(W,L)$ be True when word $W$ contains letter $L$
    (e.g.\ $\mathrm{WordContainsLetter}(\mathtt{quizzical},\mathtt z)$ is True).
    \begin{enumerate}[label=(\alph*)]
        \item Write the statement of Theorem 15 (\emph{Every English word has a
              vowel or a `y'}) in predicate logic.
        \item Now also use a unary predicate $\mathrm{Vowel}$ on letters. Write
              Theorem 15 again, more compactly.
    \end{enumerate}
\end{exercise}
\begin{solbox}

\begin{solution}
Write $\text{WordContainsLetter}(W,L)$ as $W(W,L)$.

\medskip
\noindent\textbf{(a)}
\[
\forall W\,\exists L\ \Big(\big(L=a\vee L=e\vee L=i\vee L=o\vee L=u\big)\wedge W(W,L)\Big)\ \vee\ W(W,y)
\]

\medskip
\noindent\textbf{(b)} Writing $\text{Vowel}(L)$ as $V(L)$:
\[
\forall W\,\exists L\ \Big(V(L)\wedge W(W,L)\Big)\ \vee\ W(W,y)
\]
or, with the disjunction brought inside the scope of $\exists L$:
\[
\forall W\ \Big(\exists L\ \big(V(L)\wedge W(W,L)\big)\ \vee\ W(W,y)\Big)
\]
\end{solution}
\end{solbox}

\begin{exercise}{10}
    Using the properties $\mathrm{Even}$, $\mathrm{Prime}$ and the binary
    predicate $\le$ on $\N$, write a predicate logic expression for:
    \emph{There are infinitely many odd prime numbers but only finitely many even
    prime numbers.}
\end{exercise}
\begin{solbox}

\begin{solution}
Write $\text{Even}(n)$ as $E(n)$ and $\text{Prime}(n)$ as $P(n)$. By
the domain assumption, any element that is not even is odd, so
``odd'' is $\neg E(n)$, and ``odd prime'' is $P(n)\wedge\neg E(n)$.

\medskip
\noindent\textbf{Infinitely many odd primes.} There is no largest
odd prime -- for any bound $b$, some odd prime exceeds it:
\[
R = \forall b\,\exists n\ \big(P(n)\wedge\neg E(n)\wedge b\leq n\big)
\]

\medskip
\noindent\textbf{Only finitely many even primes.} There is a
largest even prime $m$, and every even prime is at most $m$:
\[
S = \exists m\ \Big(P(m)\wedge E(m)\ \wedge\ \forall n\big(P(n)\wedge E(n)\Rightarrow n\leq m\big)\Big)
\]

\medskip
\noindent\textbf{Final statement.}
\[
R\wedge S
\]
\end{solution}
\end{solbox}

\begin{exercise}{11}
    The ternary predicate $\mathrm{Date}(d,m,y)$ on $\N^3$ is True iff
    $(d,m,y)$ is a valid Gregorian date in day-month-year format. Using
    $\mathrm{Date}$ and $<$, write an expression which is True iff
    $(d_1,m_1,y_1)$ and $(d_2,m_2,y_2)$ are both valid dates and the first comes
    chronologically before the second.
\end{exercise}
\begin{solbox}

\begin{solution}
Write $\text{Date}(d,m,y)$ as $D(d,m,y)$. $(d_1,m_1,y_1)$ comes
chronologically before $(d_2,m_2,y_2)$ exactly when
\[
y_1<y_2, \quad\text{or}\quad y_1=y_2 \wedge m_1<m_2, \quad\text{or}\quad y_1=y_2\wedge m_1=m_2\wedge d_1<d_2.
\]

So the answer is
\[
D(d_1,m_1,y_1)\wedge D(d_2,m_2,y_2)\wedge\Big((y_1<y_2)\vee\big((y_1=y_2)\wedge(m_1<m_2)\big)\vee\big((y_1=y_2)\wedge(m_1=m_2)\wedge(d_1<d_2)\big)\Big)
\]
\end{solution}
\end{solbox}

\begin{exercise}{12}
    Suppose you have: the predicate $\le$; the function
    $t:\{\text{programs}\}\times A^*\to\N$, $t(P,x)$ = time taken by program $P$
    on input $x$; the function $\mathrm{len}:A^*\to\N_0$; multiplication and
    exponentiation on positive integers; a string variable $x\in A^*$; positive
    integer variables $c,k$. Write in predicate logic:
    \emph{There exist integers $c$ and $k$ such that, for every input string $x$,
    the time taken by program $P$ on input $x$ is at most the product of $c$ and
    the $k$-th power of the length of $x$.}
\end{exercise}

\begin{exercise}{13}
    $\exists!x\,P(x)$ means ``there exists a unique $x$ such that $P(x)$
    holds''. How would you write this statement using only the tools of
    predicate logic (without ``!'')?
\end{exercise}
\begin{solbox}

\begin{solution}
\[
\exists! x\ P(x) \;\equiv\; \exists x\ \Big(P(x)\wedge\big(\forall y\ (P(y)\Rightarrow x=y)\big)\Big)
\]
\end{solution}
\end{solbox}

\begin{exercise}{14}
    \begin{enumerate}[label=(\alph*)]
        \item Prove that $\exists X\,(P(X)\wedge Q(X))$ implies
              $(\exists X\,P(X))\wedge(\exists X\,Q(X))$.
        \item Prove that $\forall X\,(P(X)\vee Q(X))$ is implied by
              $(\forall X\,P(X))\vee(\forall X\,Q(X))$ in two ways: (i) directly;
              (ii) using (a) and the relationship between existential and
              universal quantifiers.
    \end{enumerate}
\end{exercise}

